\documentclass[11pt,a4paper]{article}

\usepackage[spanish,es-noquoting]{babel}
\usepackage{fontspec}
\usepackage[margin=2.5cm]{geometry}
\usepackage{booktabs}
\usepackage{tabularx}
\usepackage{enumitem}
\usepackage{xcolor}
\usepackage[hidelinks]{hyperref}

\setlist{nosep, leftmargin=*}
\setlength{\parskip}{0.5em}
\setlength{\parindent}{0pt}

\newcommand{\hecho}{\textcolor{teal}{\textbf{Hecho}}}
\newcommand{\encurso}{\textcolor{orange!80!black}{\textbf{En curso}}}
\newcommand{\pendiente}{\textcolor{gray}{\textbf{Pendiente}}}
\newcommand{\abierta}[1]{\textcolor{orange!80!black}{\textbf{Abierta:}} #1}

\title{Avance del proyecto\\[0.3em]
  \large Armonización de focos de calor MODIS--VIIRS}
\author{Equipo Hots-Spots --- NASA Space Apps Challenge 2026}
\date{14 de noviembre de 2026}

\begin{document}
\maketitle

\section{Para qué es este documento}

Reunir en un solo lugar lo que llevamos hecho, cómo reproducirlo y qué sigue. Los
detalles de cada parte están en sus propios documentos:

\begin{tabularx}{\linewidth}{@{}lX@{}}
\toprule
Documento & Contenido \\
\midrule
\texttt{README.md} & El reto, el problema con nuestros datos y los criterios de éxito. \\
\texttt{analisis/correlacion-modis-viirs.md} & Correlación MODIS--VIIRS y constantes de conversión. \\
\texttt{analisis/serie-armonizada.md} & Regla de armonización y su validación. \\
\texttt{Documentacion/frontend/plan.tex} & Propuesta de la aplicación web. \\
\bottomrule
\end{tabularx}

\section{Estado general}

\begin{tabularx}{\linewidth}{@{}lXl@{}}
\toprule
Frente & Qué incluye & Estado \\
\midrule
Datos de prueba & CSV de FIRMS para India, 2000--2025, cargados en DuckDB & \hecho \\
Correlación & Constante por mes del año, validada con 2019--2025 & \hecho \\
Armonización & Serie 2000--2025 en escala VIIRS S-NPP, sin salto en 2012 & \hecho \\
Productos para la app & Exportar a Parquet o rejilla$\times$mes, incertidumbre & \encurso{} (equipo de datos) \\
Backend & Esqueleto FastAPI + DuckDB (solo \texttt{/salud}) & \encurso \\
Entorno & mise y Docker / Podman & \hecho \\
Frontend & Plan (fase 0) & \hecho \\
Frontend & Prototipo con datos simulados a partir de VIIRS & \pendiente \\
\bottomrule
\end{tabularx}

\section{Datos}

CSV anuales de FIRMS para \textbf{India}, en \texttt{../data de prueba/} (fuera del repo,
1{,}1\,GB). Se descargan sin clave desde
\begin{center}
\footnotesize\texttt{https://firms.modaps.eosdis.nasa.gov/data/country/<sensor>/<año>/<sensor>\_<año>\_India.csv}
\end{center}
con \texttt{<sensor>} $\in$ \{\texttt{modis} (2000--2025), \texttt{viirs-snpp} (2012--2025),
\texttt{viirs-jpss1} (2018--2025)\}.

Filtros comunes: solo vegetación (\texttt{type = 0}), MODIS con \texttt{confidence $\geq$ 30}
y VIIRS con \texttt{confidence} \texttt{n} o \texttt{h}.

\section{Armonización}

\subsection*{El problema}
VIIRS detecta entre 6 y 8 veces más focos que MODIS sobre el mismo territorio. Empalmar
MODIS hasta 2011 con VIIRS desde 2012 produce un aumento falso de $\times$5{,}4 en la serie.

\subsection*{El método}
La relación es proporcional pura (regresión log-log con exponente $b = 1{,}002$, $r = 0{,}984$),
así que basta una constante. Se usa \textbf{una constante por mes del año}, calibrada con
2012-02 a 2018-12, y se aplica \textbf{por foco} como un \texttt{peso}:

\begin{tabularx}{\linewidth}{@{}llX@{}}
\toprule
Periodo & Fuente & Peso de cada foco \\
\midrule
2000-11 a 2002-07 & Solo Terra & $k_{\mathrm{terra}}$ del mes del año \\
2002-08 a 2012-01 & Terra + Aqua & $k_{\mathrm{modis}}$ del mes del año \\
2012-02 en adelante & VIIRS S-NPP & 1 \\
\bottomrule
\end{tabularx}

Al ir por foco, los focos armonizados de cualquier área son $\sum \texttt{peso}$, sin recalcular nada.

\subsection*{Resultados}

\begin{tabularx}{\linewidth}{@{}Xr@{}}
\toprule
Indicador & Valor \\
\midrule
Razón 2012--14 / 2009--11, empalme crudo & 5{,}36 \\
Razón 2012--14 / 2009--11, serie armonizada & \textbf{0{,}96} \\
Validación 2019--2025: correlación mensual & 0{,}987 \\
Validación 2019--2025: error mediano por mes & 10{,}3\,\% \\
Validación 2019--2025: sesgo & $+2{,}7$\,\% \\
Solo Terra frente a MODIS$\times k$ (2003--2011): error mediano & 11{,}8\,\% \\
\bottomrule
\end{tabularx}

NOAA-20 y S-NPP son intercambiables ($k = 1{,}014$, $r = 0{,}999$), así que los dos VIIRS no
necesitan armonizarse entre sí.

\subsection*{Limitaciones}
\begin{itemize}
  \item Las constantes valen para India en conjunto; otra región exige recalcularlas.
  \item 2000--2002 (solo Terra) es menos preciso y la app debe marcarlo.
  \item En áreas pequeñas y en datos diarios el error relativo es mayor.
\end{itemize}

\section{Estructura del repositorio}

\begin{tabularx}{\linewidth}{@{}lX@{}}
\toprule
Carpeta & Contenido \\
\midrule
\texttt{analisis/} & Scripts SQL de DuckDB (\texttt{01}--\texttt{04}) y sus informes. La base
  \texttt{focos.duckdb} se genera localmente y no se sube. \\
\texttt{backend/} & API en FastAPI con DuckDB, gestionada con uv. \\
\texttt{frontend/} & React 19 + TypeScript + Vite + Tailwind 4. \\
\texttt{Documentacion/} & Este documento y el plan del frontend (\texttt{frontend/}). \\
\texttt{compose.yaml} & Entorno de desarrollo en contenedores. \\
\texttt{mise.toml} & Versiones de herramientas y tareas. \\
\bottomrule
\end{tabularx}

\section{Entorno de trabajo}

\subsection*{Con mise (local)}
\begin{tabularx}{\linewidth}{@{}lX@{}}
\toprule
Tarea & Qué hace \\
\midrule
\texttt{mise install} & Instala Node, pnpm, Python, uv, DuckDB y Tectonic. \\
\texttt{mise run datos:cargar} & CSV $\to$ \texttt{analisis/focos.duckdb}. \\
\texttt{mise run datos:correlacion} & Correlación y constantes por mes. \\
\texttt{mise run datos:armonizar} & Serie armonizada 2000--2025. \\
\texttt{mise run front:dev} & Servidor de Vite en \texttt{localhost:5173}. \\
\texttt{mise run back:dev} & API en \texttt{localhost:8000}. \\
\texttt{mise run docs:pdf} & Compila los \texttt{.tex} de \texttt{Documentacion/}. \\
\bottomrule
\end{tabularx}

\subsection*{Con Docker o Podman}
\begin{tabularx}{\linewidth}{@{}lX@{}}
\toprule
Tarea & Qué hace \\
\midrule
\texttt{mise run docker:up} & Construye y levanta frontend (5173) y backend (8000) con recarga. \\
\texttt{mise run docker:down} & Detiene y elimina los contenedores. \\
\bottomrule
\end{tabularx}

\begin{itemize}
  \item \texttt{frontend/Dockerfile} tiene tres etapas: \texttt{dev} (Vite), \texttt{build}
        y \texttt{prod} (nginx sirviendo el sitio estático).
  \item El backend monta \texttt{analisis/} en \texttt{/datos} en solo lectura; la ruta de la
        base está en la variable \texttt{DUCKDB\_PATH}.
  \item Los volúmenes llevan \texttt{:z} para funcionar con SELinux (Fedora).
  \item \textbf{Docker Desktop en Windows:} sin mise basta con
        \texttt{docker compose up --build} desde la raíz del repo. La recarga en caliente
        usa sondeo (\texttt{WATCH\_POLLING} en Vite y \texttt{WATCHFILES\_\allowbreak FORCE\_\allowbreak POLLING} en
        FastAPI), porque los cambios en carpetas de Windows no llegan al contenedor como
        eventos. \texttt{.gitattributes} fuerza finales de línea LF al clonar.
\end{itemize}

\section{Frontend: siguiente paso}

\subsection*{Simular la app con un CSV de VIIRS}
Para no esperar a los productos finales del equipo de datos, el prototipo se alimenta de
un \textbf{CSV de VIIRS S-NPP}. No es una aproximación arbitraria: VIIRS S-NPP \emph{es} la
escala de referencia de la armonización (peso $= 1$), así que para 2012--2025 esos focos ya son
los focos armonizados.

\begin{enumerate}
  \item Agregar los CSV de VIIRS S-NPP a rejilla (p.\,ej. $0{,}1^\circ$) por mes, con el
        mismo filtro (\texttt{type = 0}, confianza \texttt{n}/\texttt{h}).
  \item Guardar el resultado como archivo de ejemplo con el esquema del contrato de datos
        (sección ``Contrato de datos'' del plan del frontend).
  \item Construir el calendario, la frase de respuesta y la comparación con lo normal sobre
        ese archivo.
  \item Cuando llegue \texttt{focos\_armonizados}, cambiar el archivo y extender a 2000--2011,
        sin tocar los componentes.
\end{enumerate}

Lo que \emph{no} se puede simular con VIIRS: los años 2000--2011 ni la banda de
incertidumbre de MODIS. En el prototipo esas filas aparecen vacías o con datos sintéticos
marcados como tales.

\section{Decisiones abiertas}

\begin{enumerate}
  \item \abierta{¿El frontend obtiene los datos de la API FastAPI o de archivos estáticos
        (Parquet + DuckDB-WASM)? El plan del frontend proponía no tener backend, pero ya existe
        uno. Con la simulación esto se puede posponer.}
  \item \abierta{Región de la demo final: India (datos que ya tenemos) o Argentina (CONAE).}
  \item \abierta{``Normal'': mediana o promedio, de todo el periodo o de una ventana móvil.}
  \item \abierta{Idioma de la interfaz: español, o español e inglés.}
\end{enumerate}

\end{document}
